\section{Conclusion}

This work studied low-light image enhancement with a specific focus on cross-dataset generalization. All model families were trained on LOL-v2 Real and evaluated on its held-out test split, on the independent external datasets LOL-v2 Synthetic, MILLs, and ExDark, and separately on LOL-v1 as a non-independent overlap-control set. This protocol made it possible to evaluate not only reconstruction quality, but also robustness under distribution shift and category-dependent behavior on unpaired data. It also kept independent evidence separate from overlap-controlled observations, so that performance on shared LOL-v1 content could not be mistaken for external transfer.

The project implemented a complete experimental pipeline, from dataset manifest generation and deterministic preprocessing to training, validation-based model selection, multi-seed evaluation, log aggregation, and qualitative export. The final comparison was based on 624 logged experimental runs and combined full-reference metrics, no-reference metrics, diagnostic indicators, and visual failure analysis. This combination was important because low-light enhancement cannot be judged from a single score. A model may improve visibility while still producing residual darkness, clipping, color cast, over-smoothing, residual noise, halos, or hallucinated details. Looking jointly at these outputs was essential because an improved score did not necessarily correspond to faithful color, artifact-free detail, or stable exposure.

The first model family, the classical U-Net, provided a strong and useful reference. It performed well on LOL-v2 Real. Its LOL-v1 results are reported only as descriptive overlap-control evidence, rather than as independent transfer evidence. However, its behavior degraded on more shifted datasets, especially MILLs and ExDark, where color imbalance, residual darkness, and local artifacts became more visible.

The illumination-guided U-Net tested whether an explicit brightness cue was sufficient to improve generalization. The results showed that this idea was only partially effective. The illumination map helped reduce residual darkness in some cases, especially on ExDark and on darker MILLs levels, but it did not consistently improve paired reconstruction. This was one of the main lessons of the project: making the model more aware of local brightness does not automatically make the enhancement more faithful to the target or more stable in color.

FAC-ECL improved the supervision strategy by adding frequency-aware losses and exposure consistency. This family gave the safest quantitative improvement, especially on MILLs, where the test distribution differs more clearly from LOL-v2 Real while still providing paired targets. The result supports the idea that cross-dataset enhancement benefits from separating broad illumination and color behavior from local detail behavior, and from encouraging the model to remain stable under mild photometric perturbations.

ECL-FAR was introduced as an original refinement-based extension of this idea. Instead of using photometric consistency only as a regularizer, it uses the model's own sensitivity to photometric perturbations as an explicit cue for local refinement. Although ECL-FAR was implemented and tested as a first version, without an exhaustive optimization of all its components, it was already competitive with FAC-ECL and obtained the strongest structural similarity result on MILLs. In selected qualitative comparisons, it also often preserved colors better and kept dark regions cleaner than the other models.

At the same time, the results also show that ECL-FAR is not a complete solution. On LOL-v2 Synthetic, it was weaker than the other families, confirming that synthetic low-light degradation can be a difficult and partly misleading distribution shift for models trained on real low-light acquisitions. On ExDark, the no-reference and diagnostic metrics were mixed: ECL-FAR obtained the best NIQE but also showed more residual darkness and stronger RGB imbalance. Therefore, its refinement mechanism is promising, but it still needs better control of color, clipping, and extreme under-exposure.

Overall, the main conclusion of this work is that low-light enhancement should not be reduced to brightness amplification. The main challenges encountered in the project were precisely those that simple brightening cannot solve: domain shift between real and synthetic low-light formation, weak input signal in extremely dark regions, color instability across datasets, amplification of noise, halo artifacts near strong edges, and unreliable behavior on unpaired real-world scenes. The most effective models are those that constrain how enhancement is performed: how exposure is corrected, how colors are preserved, how local details are reconstructed, and how the model reacts to changes in the input distribution. FAC-ECL and ECL-FAR both support this conclusion: FAC-ECL provides the safest paired reconstruction profile, while ECL-FAR is a promising refinement-based family with the strongest observed structural result on MILLs.

Future work should study ECL-FAR more deeply. The instability map could be computed, made channel-aware, or combined with luminance, saturation, and edge information. The refinement module could be made multi-scale or separated into brightness, color, and detail corrections. The objective could also include dedicated terms for clipping control, color consistency, denoising, and edge-aware refinement. These directions would help determine whether the promising first results of ECL-FAR can be turned into a more stable and general low-light enhancement method. A systematic ablation of these choices would also clarify which components drive the gains, which interact with one another, and which cause the remaining failures.